\chapter*{How to read this guide}
\addcontentsline{toc}{chapter}{How to read this guide}
\markboth{How to read this guide}{}

Welcome! This guide explains an entire deep-learning codebase --- a TensorFlow
reimplementation of \textbf{YOLOv8} that also predicts \textbf{object outlines}
(polygons) and \textbf{how far away} each object is. It is written for someone who
is comfortable with Python but is \emph{new} to object detection. We start from the
ideas, then walk through the actual code, file by file.

\section*{Who this is for}
You will get the most out of this if you know basic Python and a little about neural
networks (``a model takes an input, has weights, and is trained with gradient
descent''). You do \emph{not} need prior knowledge of object detection, anchors,
IoU, segmentation, or YOLO. Every such term is introduced when first used.

\section*{The four kinds of callout boxes}
Throughout the guide you will see colored boxes. Each has a job:

\begin{teacher}
Explains a concept from scratch --- what it means and why it exists. Read these if a
term is new to you.
\end{teacher}

\begin{hood}
Shows the concrete mechanism in \emph{this} codebase: the function, the tensor shape,
the exact formula. Read these to connect the idea to the code.
\end{hood}

\begin{why}
A non-obvious design decision and the reasoning behind it. Several choices here differ
from ``stock'' YOLOv8 on purpose; these boxes say why.
\end{why}

\begin{pitfall}
A subtle trap --- a coordinate-order mismatch, a sentinel value, a normalization
convention --- that causes real bugs if you get it wrong.
\end{pitfall}

\section*{Conventions}
File and symbol names look like \file{data\_pipeline/mosaic.py} and \code{Mosaic.\_mosaic}.
Tensor shapes are written \code{[N, 72]}. Where coordinates matter we always say which
convention is in play --- \code{yxyx} vs \code{xyxy}, normalized \code{[0,1]} vs pixels ---
because mixing them up is the single most common source of bugs in this kind of code.

\section*{A map of the journey}
\begin{enumerate}[leftmargin=1.4em,itemsep=2pt]
  \item \textbf{Foundations} --- the object-detection ideas you need (boxes, IoU, NMS,
        anchors, feature pyramids, segmentation, the YOLO idea).
  \item \textbf{The system at a glance} --- the three model tiers and the six heads.
  \item \textbf{The data pipeline} --- how images and labels are loaded and augmented.
  \item \textbf{Polygons} --- the radial ``PolyYOLO'' representation, in depth.
  \item \textbf{The model} --- backbone, decoder, heads, and post-processing.
  \item \textbf{The losses} --- task-aligned assignment and every loss term.
  \item \textbf{Training} --- optimizer, schedules, EMA, the training loop.
  \item \textbf{Configuration, evaluation, tooling} --- how to actually run it.
\end{enumerate}

\noindent Let's begin.
